\chapter{Funkcje pierwotne i podstawowe twierdzenie rachunku całkowego}
\label{ma2:ch:primitives}

Definicję skalarnej funkcji pierwotnej, jednoznaczność
z dokładnością do stałej, tablicę funkcji pierwotnych
i warunki konieczne ich istnienia omówiono w
\atlasref{AM1-28::ma1:ch:primitives}{analizie matematycznej I}.
Z kolei podstawowe twierdzenie rachunku całkowego
i wzór Newtona--Leibniza w przypadku skalarnym
znajdują się w
\atlasref{AM1-29::ma1:thm:riemann-fundamental}{AM I, rozdział 29}.
Nie będziemy powtarzać metod obliczania całek
nieoznaczonych z tomu I.
Naszym celem jest ujęcie funkcji pierwotnej jako
rozwiązania równania różniczkowego o zadanej
pochodnej w przestrzeni Banacha oraz jako obrazu
funkcji przez ciągły operator całkowania.
W rozdziale \ref{ma2:ch:limits-calculus} wprowadziliśmy
już sumy Riemanna dla funkcji wektorowych,
a w rozdziale \ref{ma2:ch:mean-value} korzystaliśmy
ze wzoru całkowego na przyrost. Tu uzupełnimy
zależności potrzebne do konstrukcji funkcji pierwotnych;
ogólniejsze rodzaje całek wektorowych pozostawimy
do następnego rozdziału.

\section{Funkcja pierwotna krzywej}

Niech $E$ będzie rzeczywistą lub zespoloną
przestrzenią Banacha, a $J\subseteq\R$
niepustym otwartym przedziałem.
Zmienna niezależna jest zawsze rzeczywista.
W przypadku zespolonego $E$ pochodną krzywej
rozumiemy jako granicę ilorazów różnicowych
dla rzeczywistych przyrostów.

\begin{definicja}
\label{ma2:def:vector-primitive}
Funkcję różniczkowalną $F:J\to E$ nazywamy
funkcją pierwotną $f:J\to E$, gdy
$F'(t)=f(t)$ dla każdego $t\in J$.
Różniczka Frécheta krzywej $F$ w punkcie $t$
jest odwzorowaniem $DF(t):\R\to E$,
$DF(t)h=hF'(t)$.
\end{definicja}

Ostatnia równość wynika z definicji różniczki:
\[
 \frac{\|F(t+h)-F(t)-hF'(t)\|}{|h|}
 =\left\|\frac{F(t+h)-F(t)}h-F'(t)\right\|
 \longrightarrow0.
\]
Odwrotnie, dla różniczki $A:\R\to E$ mamy
$Ah=hA(1)$, więc istnienie różniczki pociąga
istnienie pochodnej $F'(t)=A(1)$.
W ten sposób utożsamiamy $\mathcal L(\R,E)$
z $E$ przez izometryczny izomorfizm $A\mapsto A(1)$
(w przypadku zespolonym jest to przestrzeń
odwzorowań rzeczywiście liniowych).

\begin{stwierdzenie}
\label{ma2:prop:primitive-uniqueness}
Jeżeli $F,G:J\to E$ są funkcjami pierwotnymi
tej samej funkcji $f$, to $F-G$ jest stała.
Jeżeli istnieje jedna funkcja pierwotna $F$,
to wszystkie mają postać $F+c$, $c\in E$.
Dla $t_0\in J$ i $v\in E$ istnieje wtedy
dokładnie jedna funkcja pierwotna przyjmująca
w $t_0$ wartość $v$.
\end{stwierdzenie}
\begin{proof}
Dla $s<t$ w $J$ stosujemy oszacowanie przyrostu
\ref{ma2:thm:mean-value-inequality} do $H=F-G$
na $[s,t]$, z $M=0$. Dostajemy
$\|H(t)-H(s)\|\le0$, więc $H(t)=H(s)$.
Dodanie stałej nie zmienia pochodnej.
Wartość w $t_0$ wymusza
$c=v-F(t_0)$, a ten wybór rzeczywiście daje
żądaną funkcję.
\end{proof}

Jeżeli dziedziną jest dowolny otwarty podzbiór
$\R$, to jego składowe spójności są otwartymi
przedziałami, a stałą można wybierać niezależnie
na każdej składowej.
W szczególności funkcja pierwotna na
$\R\setminus\{0\}$ nie musi różnić się od innej
funkcji pierwotnej o jedną wspólną stałą.

\section{Całka po podprzedziale}

Przypomnijmy, że całkowalność Riemanna
$f:[a,b]\to E$ oznacza zbieżność w normie
wszystkich sum
\[
 R(f,P)=\sum_{j=1}^m(t_j-t_{j-1})f(\xi_j),
 \qquad \xi_j\in[t_{j-1},t_j],
\]
przy średnicy podziału dążącej do zera, jednostajnie
względem wyboru punktów pośrednich.
Ograniczoność funkcji całkowalnej wykazano
w rozdziale \ref{ma2:ch:limits-calculus}.
Dla funkcji ciągłych istnienie całki zapewnia
zupełność przestrzeni Banacha (por.
\ref{ma2:lem:ftc-banach} i odsyłacz do AF 15A
w rozdziale \ref{ma2:ch:limits-calculus}).

\begin{lemat}[Kryterium Cauchy'ego dla sum Riemanna]
\label{ma2:lem:primitive-cauchy}
Funkcja $f:[a,b]\to E$ jest całkowalna w sensie
Riemanna wtedy i tylko wtedy, gdy dla każdego
$\varepsilon>0$ istnieje $\delta>0$ takie, że
\[
 \|R(f,P)-R(f,T)\|<\varepsilon
\]
dla wszystkich podziałów z punktami pośrednimi
$P,T$ o średnicach mniejszych od $\delta$.
\end{lemat}
\begin{proof}
Jeżeli całka istnieje, wybieramy $\delta$
tak, aby obie sumy różniły się od całki
o mniej niż $\varepsilon/2$, i stosujemy
nierówność trójkąta.

Odwrotnie, wybierzmy ciąg podziałów $P_m$
o średnicach dążących do zera, z dowolnymi
punktami pośrednimi. Warunek Cauchy'ego
daje zbieżność $R(f,P_m)\to I\in E$
dzięki zupełności $E$.
Dla $\varepsilon>0$ bierzemy $\delta$
z warunku Cauchy'ego dla $\varepsilon/2$.
Wybieramy $m$ z $|P_m|<\delta$ i
$\|R(f,P_m)-I\|<\varepsilon/2$.
Wtedy dla każdego $P$ o $|P|<\delta$
\[
 \|R(f,P)-I\|
 \le\|R(f,P)-R(f,P_m)\|+\|R(f,P_m)-I\|
 <\varepsilon.
\]
To jest definicja całkowalności.
\end{proof}

\begin{twierdzenie}[Ograniczenie i addytywność]
\label{ma2:thm:primitive-subinterval}
Jeżeli $f:[a,b]\to E$ jest całkowalna w sensie
Riemanna, to jej ograniczenie do każdego
podprzedziału domkniętego jest całkowalne.
Dla $a\le u\le v\le w\le b$ zachodzi
\[
 \int_u^w f(t)\,\mathrm dt
 =\int_u^v f(t)\,\mathrm dt+\int_v^w f(t)\,\mathrm dt.
\]
Jeżeli $a\le u\le v\le b$, to
\[
 \left\|\int_u^v f(t)\,\mathrm dt\right\|
 \le(v-u)\sup_{t\in[u,v]}\|f(t)\|.
\]
Całce na przedziale jednopunktowym przypisujemy
wartość zero.
\end{twierdzenie}
\begin{proof}
Niech $u<v$. Dla dwóch podziałów z punktami
pośrednimi $P,T$ przedziału $[u,v]$
uzupełnijmy oba tymi samymi podziałami z punktami
pośrednimi na $[a,u]$ i $[v,b]$ (pomijamy
przedziały jednopunktowe). Jeśli wszystkie
średnice są mniejsze od $\delta$ z kryterium
Cauchy'ego dla $f$, to złożone podziały
$P',T'$ całego $[a,b]$ też mają średnice
mniejsze od $\delta$. Wkłady z przedziałów
zewnętrznych znoszą się w różnicy sum, więc
\[
 \|R(f|_{[u,v]},P)-R(f|_{[u,v]},T)\|
 =\|R(f,P')-R(f,T')\|<\varepsilon.
\]
Lemat \ref{ma2:lem:primitive-cauchy} daje
całkowalność ograniczenia.

Dla $u<v<w$ wybieramy ciągi podziałów
$[u,v]$ i $[v,w]$ o średnicach dążących do zera.
Ich złożenie jest podziałem $[u,w]$
o średnicy równej maksimum obu średnic.
Suma Riemanna dla złożenia jest dokładnie
sumą dwóch sum lokalnych. Przejście do granicy
daje addytywność. Gdy $u=v$ lub $v=w$,
równość wynika z konwencji całki zerowej.

Dla dowolnego podziału $[u,v]$ nierówność
trójkąta daje
\[
 \|R(f,P)\|\le
 \sum_j(t_j-t_{j-1})\|f(\xi_j)\|
 \le(v-u)\sup_{t\in[u,v]}\|f(t)\|.
\]
Przechodzimy do granicy, korzystając
z ciągłości normy.
\end{proof}

Dla $u>v$ definiujemy całkę z odwróconymi granicami
jako $-\int_v^u f(t)\,\mathrm dt$.
Dzięki udowodnionej addytywności dla uporządkowanych
punktów równość
$\int_u^w f=\int_u^v f+\int_v^w f$
zachodzi przy dowolnej kolejności końców:
wystarczy zapisać każdą całkę jako różnicę
dwóch całek z ustalonego lewego końca $a$.
Liniowość wynika z liniowości sum dla wspólnych
podziałów i przejścia do granicy.
Jeżeli $L:E\to F$ jest ciągłym odwzorowaniem
liniowym między przestrzeniami Banacha (nad
tym samym ciałem), to
\[
 \int_a^b L(f(t))\,\mathrm dt
 =L\left(\int_a^b f(t)\,\mathrm dt\right).
\]
Istotnie, $R(L\circ f,P)=L(R(f,P))$ dla każdego
podziału, a ciągłość $L$ pozwala przejść do granicy.
To rozumowanie wykazuje zarazem całkowalność
$L\circ f$.

\section{Całka jako funkcja górnej granicy}

\begin{twierdzenie}[Podstawowe twierdzenie rachunku całkowego]
\label{ma2:thm:primitive-ftc}
Niech $f:[a,b]\to E$ będzie całkowalna
w sensie Riemanna, $t_0\in[a,b]$ i $v\in E$.
Funkcja
\[
 F(x)=v+\int_{t_0}^x f(t)\,\mathrm dt
\]
jest Lipschitzowska na $[a,b]$ ze stałą
$\sup_{t\in[a,b]}\|f(t)\|$.
W każdym punkcie ciągłości $f$ we wnętrzu
przedziału zachodzi $F'(x)=f(x)$.
Na końcach analogiczna teza dotyczy pochodnych
jednostronnych przy odpowiedniej ciągłości $f$.
\end{twierdzenie}
\begin{proof}
Z addytywności i oszacowania w twierdzeniu
\ref{ma2:thm:primitive-subinterval} mamy
\[
 \|F(y)-F(x)\|\le |y-x|\sup_{t\in[a,b]}\|f(t)\|.
\]
Niech $x$ będzie punktem ciągłości $f$.
Dla $h\ne0$, $x+h\in[a,b]$, odejmujemy
całkę funkcji stałej $f(x)$, uzyskując
\[
 \frac{F(x+h)-F(x)}h-f(x)
 =\frac1h\int_x^{x+h}(f(t)-f(x))\,\mathrm dt.
\]
Norma prawej strony nie przekracza
\[
 \sup_{\substack{t\in[a,b]\\|t-x|\le |h|}}
 \|f(t)-f(x)\|.
\]
Oszacowanie jest poprawne również dla $h<0$,
ponieważ odwrócenie granic zmienia znak,
a dzielenie przez $h$ zmienia normę przez
czynnik $1/|h|$. Dla dowolnego $\varepsilon>0$
ciągłość $f$ w $x$ daje $\delta>0$ tak, że
$\|f(t)-f(x)\|<\varepsilon/2$ dla
$|t-x|<\delta$. Jeżeli $|h|<\delta$,
supremum nie przekracza $\varepsilon/2$.
To dowodzi tezy, również przy dopuszczalnych
przyrostach jednostronnych na końcach.
\end{proof}

\begin{wniosek}
\label{ma2:cor:primitive-exists}
Każda ciągła funkcja $f:J\to E$ ma funkcję
pierwotną klasy $C^1$, a dla danych
$t_0\in J$, $v\in E$ jedyną taką funkcją
spełniającą $F(t_0)=v$ jest
\[
 F(x)=v+\int_{t_0}^x f(t)\,\mathrm dt.
\]
\end{wniosek}
\begin{proof}
Każdy domknięty odcinek o końcach w $J$
jest zawarty w $J$. Ograniczenie $f$
do niego jest ciągłe, więc całka istnieje
w przestrzeni Banacha.
Dla dowolnego $x\in J$ wybieramy domknięty
odcinek w $J$ zawierający $x$ we wnętrzu
oraz $t_0$. Poprzednie twierdzenie daje
$F'(x)=f(x)$. Ciągłość pochodnej wynika
z ciągłości $f$, a jednoznaczność ze
stwierdzenia \ref{ma2:prop:primitive-uniqueness}.
\end{proof}

\begin{przyklad}[Ciągłość jest warunkiem dostatecznym,
nie koniecznym]
\label{ma2:ex:primitive-jump}
Niech $e\ne0$ w $E$. Dla
\[
 f(t)=
 \begin{cases}0,&t<0,\\ e,&t\ge0,\end{cases}
 \quad -1\le t\le1,
\]
funkcja $f$ jest całkowalna: jest iloczynem
funkcji skalarnej o jednej nieciągłości
i ustalonego wektora (sumy Riemanna są
iloczynami odpowiednich sum skalarnych i $e$).
Jej całka od zera wynosi $F(t)=\max(t,0)e$;
wartość w pojedynczym punkcie nie wpływa na
całkę skalarną
(\atlasref{AM1-29::ma1:cor:riemann-finite-changes}{zmiany skończone}).
W zerze funkcja $F$ ma różne pochodne
jednostronne, odpowiednio $0$ oraz $e$.
Co więcej, $f$ nie ma funkcji pierwotnej
na $(-1,1)$: hipotetyczna funkcja $G$
byłaby stała na $(-1,0)$ i równa $te+c$
na $(0,1)$, z jednoznaczności funkcji pierwotnych.
Ciągłość w zerze wymusza tę samą stałą po obu
stronach, a różniczkowalność w zerze jest
wtedy niemożliwa.
\end{przyklad}

% Następny blok: Newton--Leibniz dla całkowalnej
% pochodnej, operator całkowania, całkowanie przez
% części i podstawienie, przykłady w przestrzeniach Banacha.
